\documentclass[11pt]{article}
\usepackage{amsmath,amssymb,amsthm}
\usepackage{booktabs}
\usepackage[margin=1in]{geometry}
\newtheorem{theorem}{Theorem}
\newtheorem{lemma}[theorem]{Lemma}
\title{Problem 797: Cyclogenic Polynomials}
\author{Project Euler Solutions}
\date{}
\begin{document}
\maketitle
\section{Problem Statement}
A monic polynomial $p(x)$ is $n$-cyclogenic if $p(x)q(x) = x^n - 1$ for some monic $q(x) \in \mathbb{Z}[x]$, with $n$ minimal. $P_n(x)$ = sum of all $n$-cyclogenic polynomials. $Q_N(x) = \sum_{n=1}^N P_n(x)$. Given $Q_{10}(2) = 5598$. Find $Q_{10^7}(2) \bmod 10^9+7$.
\section{Analysis}
An $n$-cyclogenic polynomial divides $x^n - 1$ in $\mathbb{Z}[x]$ and does not divide $x^m - 1$ for any $m < n$. The divisors of $x^n - 1$ are products of cyclotomic polynomials $\Phi_d(x)$ for $d | n$.

A monic divisor of $x^n-1$ is a product $\prod_{d \in S} \Phi_d(x)$ for some subset $S \subseteq \{d : d|n\}$. It is $n$-cyclogenic if $n$ is the minimal such value, meaning $S$ contains at least one $d$ for which $n$ is the smallest multiple.

Equivalently, the polynomial is $n$-cyclogenic iff $\text{lcm}(S) = n$, where $S$ is the set of cyclotomic indices used.

$P_n(x)$ is the sum over all 
\section{Answer}

\[
\boxed{47722272}
\]

\end{document}
